AWS PrivateLink endpoints have a fixed IP address in each Availability
Zone where they are deployed, for the life of the endpoint (until it is
deleted). These IP addresses are attached to Elastic Network Interfaces
(ENIs). AWS recommends using DNS to resolve the IP addresses for
endpoints so that downstream applications use the latest IP addresses
when ENIs are added to new AZs, or deleted over time.

In this section, you will create a forwarding rule to send DNS
resolution requests from a simulated on-premises environment to a Route
53 Private Hosted Zone. This section leverages the infrastructure
deployed by CloudFormation in the Prepare the environment section.

\vspace{0.5em}\noindent\textbf{Create DNS Alias Records for the Interface
endpoint}\par

\begin{enumerate}
\def\labelenumi{\arabic{enumi}.}
\tightlist
\item
  Navigate to the
  \href{https://us-east-1.console.aws.amazon.com/route53/v2/hostedzones?region=us-east-1\#}{Route
  53 management console} (Hosted Zones section). The CloudFormation
  template you deployed in the Prepare the environment section created
  this Private Hosted Zone. Click on the name of the Private Hosted
  Zone, s3.us-east-1.amazonaws.com:
\end{enumerate}

\begin{figure}
\centering
\pandocbounded{\includegraphics[keepaspectratio,alt={hosted zone}]{5-Workshop/5.4-S3-onprem/hosted-zone.png}}
\caption{hosted zone}
\end{figure}

\begin{enumerate}
\def\labelenumi{\arabic{enumi}.}
\setcounter{enumi}{1}
\tightlist
\item
  Create a new record in the Private Hosted Zone:
\end{enumerate}

\begin{figure}
\centering
\pandocbounded{\includegraphics[keepaspectratio,alt={Create record}]{5-Workshop/5.4-S3-onprem/create-record1.png}}
\caption{Create record}
\end{figure}

\begin{itemize}
\tightlist
\item
  Record name and record type keep default options
\item
  Alias Button: Click to enable
\item
  Route traffic to: Alias to VPC Endpoint
\item
  Region: US East (N. Virginia) {[}us-east-1{]}
\item
  Choose endpoint: Paste the Regional VPC Endpoint DNS name from your
  text editor (you saved when doing section 4.3)
\end{itemize}

\begin{figure}
\centering
\pandocbounded{\includegraphics[keepaspectratio,alt={record1}]{5-Workshop/5.4-S3-onprem/record1.png}}
\caption{record1}
\end{figure}

\begin{enumerate}
\def\labelenumi{\arabic{enumi}.}
\setcounter{enumi}{2}
\tightlist
\item
  Click Add another record, and add a second record using the following
  values. Click Create records when finished to create both records.
\end{enumerate}

\begin{itemize}
\tightlist
\item
  Record name: *.
\item
  Record type: keep default value (type A)
\item
  Alias Button: Click to enable
\item
  Route traffic to: Alias to VPC Endpoint
\item
  Region: US East (N. Virginia) {[}us-east-1{]}
\item
  Choose endpoint: Paste the Regional VPC Endpoint DNS name from your
  text editor
\end{itemize}

\begin{figure}
\centering
\pandocbounded{\includegraphics[keepaspectratio,alt={record 2}]{5-Workshop/5.4-S3-onprem/record2.png}}
\caption{record 2}
\end{figure}

The new records appear in the Route 53 console:

\begin{figure}
\centering
\pandocbounded{\includegraphics[keepaspectratio,alt={result}]{5-Workshop/5.4-S3-onprem/result.png}}
\caption{result}
\end{figure}

\vspace{0.5em}\noindent\textbf{Create a Resolver Forwarding
Rule}\par

Route 53 Resolver Forwarding Rules allow you to forward DNS queries from
your VPC to other sources for name resolution. Outside of a workshop
environment, you might use this feature to forward DNS queries from your
VPC to DNS servers running on-premises. In this section, you will
simulate an on-premises conditional forwarder by creating a forwarding
rule that forwards DNS queries for Amazon S3 to a Private Hosted Zone
running in ``VPC Cloud'' in-order to resolve the PrivateLink interface
endpoint regional DNS name.

\begin{enumerate}
\def\labelenumi{\arabic{enumi}.}
\tightlist
\item
  From the Route 53 management console, click \textbf{Inbound endpoints}
  on the left side bar
\item
  In the Inbound endpoints console, click the ID of the inbound endpoint
\end{enumerate}

\begin{figure}
\centering
\pandocbounded{\includegraphics[keepaspectratio,alt={Inbound endpoint}]{5-Workshop/5.4-S3-onprem/route53-1.png}}
\caption{Inbound endpoint}
\end{figure}

\begin{enumerate}
\def\labelenumi{\arabic{enumi}.}
\setcounter{enumi}{2}
\tightlist
\item
  Copy the two IP addresses listed to your text editor
\end{enumerate}

\begin{figure}
\centering
\pandocbounded{\includegraphics[keepaspectratio,alt={Ip addresses}]{5-Workshop/5.4-S3-onprem/route53-2.png}}
\caption{Ip addresses}
\end{figure}

\begin{enumerate}
\def\labelenumi{\arabic{enumi}.}
\setcounter{enumi}{3}
\tightlist
\item
  From the Route 53 menu, choose \textbf{Resolver} \textgreater{}
  \textbf{Rules}, and click \textbf{Create rule}:
\end{enumerate}

\begin{figure}
\centering
\pandocbounded{\includegraphics[keepaspectratio,alt={Ip addresses}]{5-Workshop/5.4-S3-onprem/route53-3.png}}
\caption{Ip addresses}
\end{figure}

\begin{enumerate}
\def\labelenumi{\arabic{enumi}.}
\setcounter{enumi}{4}
\tightlist
\item
  In the Create rule console:
\end{enumerate}

\begin{itemize}
\tightlist
\item
  Name: myS3Rule
\item
  Rule type: Forward
\item
  Domain name: s3.us-east-1.amazonaws.com
\end{itemize}

\begin{figure}
\centering
\pandocbounded{\includegraphics[keepaspectratio,alt={create rule}]{5-Workshop/5.4-S3-onprem/route53-4.png}}
\caption{create rule}
\end{figure}

\begin{itemize}
\tightlist
\item
  VPC: VPC On-prem
\item
  Outbound endpoint: VPCOnpremOutboundEndpoint
\end{itemize}

\begin{figure}
\centering
\pandocbounded{\includegraphics[keepaspectratio,alt={create rule}]{5-Workshop/5.4-S3-onprem/route53-5.png}}
\caption{create rule}
\end{figure}

\begin{itemize}
\tightlist
\item
  Target IP Addresses: Enter both IP addresses from your text editor
  (inbound endpoint addresses) and then click Submit
\end{itemize}

\pandocbounded{\includegraphics[keepaspectratio,alt={create rule}]{5-Workshop/5.4-S3-onprem/route53-6.png}}
You have successfully created resolver forwarding rule.

\begin{figure}
\centering
\pandocbounded{\includegraphics[keepaspectratio,alt={create rule}]{5-Workshop/5.4-S3-onprem/route53-7.png}}
\caption{create rule}
\end{figure}

\vspace{0.5em}\noindent\textbf{Test the on-premises DNS
Simulation}\par

\begin{enumerate}
\def\labelenumi{\arabic{enumi}.}
\tightlist
\item
  Connect to \textbf{Test-Interface-Endpoint EC2 instance} with
  \textbf{Session manager}
\end{enumerate}

\begin{figure}
\centering
\pandocbounded{\includegraphics[keepaspectratio,alt={create rule}]{5-Workshop/5.4-S3-onprem/test1.png}}
\caption{create rule}
\end{figure}

\begin{enumerate}
\def\labelenumi{\arabic{enumi}.}
\setcounter{enumi}{1}
\tightlist
\item
  Test DNS resolution. The dig command will return the IP addresses
  assigned to the VPC Interface endpoint running in VPC Cloud (your IP's
  will be different): dig +short s3.us-east-1.amazonaws.com
\end{enumerate}

The IP addresses returned are the VPC endpoint IP addresses, NOT the
Resolver IP addresses you pasted from your text editor. The IP addresses
of the Resolver endpoint and the VPC endpoint look similar because they
are all from the VPC Cloud CIDR block.

\begin{figure}
\centering
\pandocbounded{\includegraphics[keepaspectratio,alt={create rule}]{5-Workshop/5.4-S3-onprem/dig.png}}
\caption{create rule}
\end{figure}

\begin{enumerate}
\def\labelenumi{\arabic{enumi}.}
\setcounter{enumi}{2}
\tightlist
\item
  Navigate to the VPC menu (Endpoints section), select the S3 Interface
  endpoint. Click the Subnets tab and verify that the IP addresses
  returned by Dig match the VPC endpoint:
\end{enumerate}

\begin{figure}
\centering
\pandocbounded{\includegraphics[keepaspectratio,alt={create rule}]{5-Workshop/5.4-S3-onprem/subnet.png}}
\caption{create rule}
\end{figure}

\begin{enumerate}
\def\labelenumi{\arabic{enumi}.}
\setcounter{enumi}{3}
\tightlist
\item
  Return to your shell and use the AWS CLI to test listing your S3
  buckets:
\end{enumerate}

\begin{verbatim}
aws s3 ls --endpoint-url https://s3.us-east-1.amazonaws.com
\end{verbatim}

\begin{figure}
\centering
\pandocbounded{\includegraphics[keepaspectratio,alt={create rule}]{5-Workshop/5.4-S3-onprem/endpoint.png}}
\caption{create rule}
\end{figure}

\begin{enumerate}
\def\labelenumi{\arabic{enumi}.}
\setcounter{enumi}{4}
\tightlist
\item
  Terminate your Session Manager session:
\end{enumerate}

\begin{figure}
\centering
\pandocbounded{\includegraphics[keepaspectratio,alt={create rule}]{5-Workshop/5.4-S3-onprem/terminal.png}}
\caption{create rule}
\end{figure}

In this section you created an Interface endpoint for Amazon S3. This
endpoint can be reached from on-premises through Site-to-Site VPN or AWS
Direct Connect. Route 53 Resolver outbound endpoints simulated
forwarding DNS requests from on-premises to a Private Hosted Zone
running the cloud. Route 53 inbound Endpoints recieved the resolution
request and returned a response containing the IP addresses of the VPC
interface endpoint. Using DNS to resolve the endpoint IP addresses
provides high availability in-case of an Availability Zone outage.
